\chapter{Literature Review}
\label{ch:review}

This chapter summarises the evidence base that drove the design inversion from
Version 1.0. A full critical review of fifty sources is provided in the companion
document \textit{Statistical and Machine-Learning Approaches to Real-Time Sepsis
Onset Prediction on MIMIC-IV: An Extended Critical Literature Review} (Version 2.0);
this chapter extracts the seven evidence strands directly load-bearing for the
present study's design choices.

%===========================================================================
\section{The Outcome Is Not One Object: Label Sensitivity}
%===========================================================================

The central premise of Version 1.0 --- that a model can be assessed on
``the MIMIC-IV sepsis cohort'' --- is undermined by two studies.

Johnson et al.\ \parencite{johnson2018comparative} applied five retrospective
sepsis identification algorithms to 11{,}791 qualifying ICU admissions from
MIMIC and found only acceptable inter-algorithm agreement (Cronbach's $\alpha$
= 0.40--0.62). The implication is not that any single algorithm is wrong, but
that ``the MIMIC sepsis cohort'' is a family of related but non-identical objects,
each constructed by a specific reading of the Sepsis-3 definition.

Cohen et al.\ \parencite{cohen2024subtle} made the consequence for prediction
modelling explicit. Holding MIMIC-III fixed and varying only the reading of
Sepsis-3 onset across three defensible interpretations, the identified cohort
ranged from 867 to 2{,}178 admissions. Across tree-based, deep-learning, and
survival-analysis models: holding the definition fixed, the best method gained
1--5\% AUROC; holding the method fixed and varying the definition swung
performance by 0--6\% AUROC. The variance attributable to the label and to the
model class are, in Cohen et al.'s data, of comparable magnitude. This finding
is the single most load-bearing source in the present design: it mandates that
label operationalisation be treated as a primary object of study, not a fixed
parameter.

Dutta et al.\ \parencite{dutta2026performance} confirm that this is not a
MIMIC artefact. A locally trained gradient-boosted model, predicting every fifteen
minutes across nine acute-care hospitals on 198{,}494 encounters, was evaluated
against Sepsis-3, SEP-1, and the CDC Adult Sepsis Event definition. Performance
varied substantially by definition. The label-sensitivity finding is therefore a
property of the problem: it reproduces prospectively, at scale, with a commercial
model.

%===========================================================================
\section{The Label Is Constituted by Clinical Action}
%===========================================================================

The Sepsis-3 label is operationalised as the combination of a qualifying
antibiotic order and a qualifying culture draw within specified windows, plus an
acute SOFA rise. Kamran et al.\ \parencite{kamran2024evaluation} observed that
clinicians frequently recognise and begin treating sepsis before the formal
criteria are met. A model trained to predict the criteria-meeting time therefore
partly learns to predict when the clinician will act. Kamran et al.\ evaluated
performance restricted to the pre-treatment window and showed that apparent AUROC
is substantially inflated without that restriction. The model had learned to
encode clinician suspicion, not independent physiological signal.

Lauritsen et al.\ \parencite{lauritsen2021framing} formalised the related
concept of \textit{framing}: the choice of observation window, prediction horizon,
window shift, and event trigger. They showed that framing changes discrimination
and calibration, and can reverse the sign of an estimated covariate effect ---
their example is SpO$_2$. This implies that a hazard ratio or odds ratio whose
sign is an artefact of framing is not a clinically interpretable quantity.

%===========================================================================
\section{The Validation Design Must Be Temporal}
%===========================================================================

Guo et al.\ \parencite{guo2022temporal} partitioned MIMIC-IV by
\texttt{anchor\_year\_group} and evaluated four prediction tasks under temporal
shift. Sepsis was the worst-affected task; domain-generalisation and
domain-adaptation algorithms did not outperform plain empirical risk minimisation.
A random patient-level split therefore inflates apparent performance more for
sepsis than for any neighbouring outcome. This finding mandates a temporal split
as the primary internal validation strategy, with a random split retained only to
quantify the optimism introduced by the conventional design.

%===========================================================================
\section{AUROC Cannot Detect the Characteristic Failure Mode}
%===========================================================================

Wang et al.\ \parencite{wang2025srpm} reviewed ninety-one real-time sepsis
prediction models. Median internal AUROC was 0.811, external AUROC was 0.783:
a difference of 0.028, not statistically significant. Median internal Utility
Score was $+0.381$, external was $-0.164$: the models crossed from clinically
useful to categorically harmful. AUROC survived the transition because it
measures discrimination, and the relative ordering of patients by risk persisted
even as the absolute scores became uncalibrated. The Utility Score failed because
it measures whether the alert arrives before the clinical window --- and at external
sites, the alert arrived too late, too often, and at unacceptable false-alert rates.
Reporting AUROC as the primary metric is therefore not merely conservative; it
actively conceals the failure mode.

%===========================================================================
\section{The Survival Model Is Available in a Form the Data Actually Takes}
%===========================================================================

Version 1.0 framed the survival approach as a Cox proportional-hazards model.
The Schoenfeld residual test rejected proportional hazards globally ($p < 10^{-16}$),
as was anticipated from Kasal et al.\ \parencite{kasal2004cox}. The appropriate
response to a failed PH test in an hourly, low-prevalence ICU regime is provided
by D'Agostino et al.\ \parencite{dagostino1990pooled}: observations over multiple
intervals are pooled and a logistic regression relates risk factors to event
occurrence within each interval. The pooled logistic regression is asymptotically
equivalent to time-dependent-covariate Cox regression for short intervals and small
per-interval event probabilities. An hourly ICU regime with per-hour onset
probability below 1\% satisfies both conditions. The dynamic discrete-time
competing-risks hazard model fitted in the present study is therefore not an
abandonment of the survival framework but its implementation in the form the data
actually takes.

Resche-Rigon et al.\ \parencite{rescherigon2006competing} and Heyard et al.\
\parencite{heyard2019dynamic, heyard2020evaluation} provide the corresponding
competing-risks extension: because discharge alive and death without sepsis are
informationally relevant events that preclude later sepsis onset, treating them
as ordinary censoring overstates the onset incidence and biases model estimation.
The multinomial cause-specific formulation accommodates this without additional
assumptions.

%===========================================================================
\section{NEWS2 Is the Appropriate Rule-Based Headline Comparator}
%===========================================================================

Version 1.0 benchmarked against SIRS, qSOFA, and NEWS. Evans et al.\
\parencite{evans2021ssc} issue a strong recommendation against using qSOFA as a
sole screening tool, on grounds of poor sensitivity. Mellhammar et al.\
\parencite{mellhammar2020news2} fitted a LASSO-derived statistical score and
failed to beat NEWS2 on the same task. The governing clinical guideline has
therefore already retired qSOFA as a screening comparator, and NEWS2 is the
minimum meaningful rule-based bar. The present study retains SIRS and qSOFA for
descriptive completeness but designates NEWS2 as the headline comparator.

%===========================================================================
\section{Equity Analysis: Subgroup Performance and Label Sensitivity Together}
%===========================================================================

Wang, Li, Naidech and Luo \parencite{wang2022equity} document significant
variation in sepsis model performance across racial, ethnic, language, and
insurance subgroups. Critically, they do not evaluate whether the label itself
assigns sepsis status differentially across those groups. The present study adds
this conjunction: for each equity subgroup, it reports both the variation in who
is identified as septic under each label variant and the variation in model
performance. This has not been reported for any sepsis model.

%===========================================================================
\section{Summary: Design Decisions and Their Evidence Warrants}
%===========================================================================

Table~\ref{tab:design_warrant} maps each major design decision in the present
study to the evidence source that mandates it.

\begin{table}[H]
  \centering
  \small
  \begin{tabularx}{\textwidth}{p{4cm} X p{4cm}}
    \toprule
    \textbf{Design decision} & \textbf{Rationale} & \textbf{Source} \\
    \midrule
    Three pre-specified label variants & Label spread comparable to model-class
      spread & Cohen et al.\ \parencite{cohen2024subtle} \\
    Label sensitivity as primary result & Not a sensitivity analysis: the variance
      is load-bearing & Cohen et al., Dutta et al.\ \parencite{dutta2026performance} \\
    Treatment-anchored evaluation & Label constituted by clinical action &
      Kamran et al.\ \parencite{kamran2024evaluation} \\
    Temporal split primary & Sepsis worst-affected task under temporal shift &
      Guo et al.\ \parencite{guo2022temporal} \\
    Utility Score as primary metric & AUROC conceals the clinical failure mode &
      Wang et al.\ \parencite{wang2025srpm} \\
    Calibration as co-primary & Achilles heel of predictive analytics &
      Van Calster et al.\ \parencite{vancalster2019calibration} \\
    Discrete-time competing-risks model & Equivalent to Cox; no PH assumption; hourly
      regime & D'Agostino et al.\ \parencite{dagostino1990pooled}; Heyard et al.\
      \parencite{heyard2019dynamic} \\
    NEWS2 as headline comparator & qSOFA not recommended as sole screener &
      Evans et al.\ \parencite{evans2021ssc} \\
    GBT as mandatory ML comparator & Cross-study comparison uninterpretable
      under different labels & Cohen et al.\ \parencite{cohen2024subtle} \\
    Subgroup label sensitivity & Novel conjunction; equity literature does
      not evaluate label-by-subgroup & Wang et al.\ \parencite{wang2022equity} \\
    \bottomrule
  \end{tabularx}
  \caption{Mapping of design decisions to their evidence warrants.}
  \label{tab:design_warrant}
\end{table}
